\section{Conclusions}
\label{sec:conclusions}

The MOSAIC pipeline demonstrates that per-patient $\Delta Z$ vectors derived from a
matched tumour-normal multi-omics cohort can support biologically informative patient
stratification.
Three reproducible transformation subtypes were identified from the CPTAC-LUAD dataset,
each characterised by a distinct transcriptional reprogramming signature.
Subtype~0 shows relatively suppressed cell-cycle activation and is significantly
enriched for \textit{EGFR} mutations.
Subtype~1 is proliferation-dominant, with strong E2F and G2-M upregulation.
Subtype~2 is characterised by interferon-driven immune activation combined with a
proliferative component.
The \textit{EGFR} enrichment in subtype~0 gains additional credibility from the
simultaneous depletion of ALK fusions in the same subtype: because \textit{EGFR}
mutations and ALK rearrangements are mutually exclusive driver events, observing
both associations together strengthens confidence in the subtype structure rather than
attributing either result to chance alone.

The geometric separation between subtypes is modest (mean silhouette $\approx 0.18$),
suggesting a partially continuous biological manifold rather than fully discrete
clusters.
No significant overall survival difference was detected, which reflects the limited
statistical power of the study ($n = 102$) rather than necessarily indicating clinical
irrelevance.
Replication in an independent LUAD cohort with matched normal profiling would be the
most informative next step.

\bigskip
\noindent\textbf{Code availability.}
The full analysis pipeline is publicly available at
\url{https://github.com/ajanowiak/MOSAIC}.
